\section{GIỚI THIỆU VÀ TÓM TẮT NHIỆM VỤ XÂY DỰNG HỒ SƠ
}
\subsection{Các căn cứ pháp lý để thực hiện dự án}
Dự án \gls{geoai} được xây dựng dựa trên các căn cứ pháp lý sau:
\begin{itemize}
    \item Luật Công nghệ thông tin số 67/2006/QH11 ngày 29/6/2006;
    \item Luật Đầu tư công số 58/2024/QH15 ngày 29 tháng 11 năm 2024;
    \item Nghị định 73/2019/\gls{nd-cp} ngày 05/9/2019 của Chính phủ về quản lý đầu tư ứng dụng công nghệ thông tin sử dụng nguồn vốn ngân sách nhà nước;
    \item Nghị định số 82/2024/\gls{nd-cp} ngày 10/7/2024 của Chính phủ về sửa đổi, bổ sung một số điều của Nghị định số 73/2019/\gls{nd-cp};
    \item Quyết định số 1688/\gls{qd-btttt} ngày 11/10/2019 của Bộ Thông tin và Truyền thông về định mức chi phí quản lý dự án và chi phí tư vấn đầu tư ứng dụng \gls{cntt};
    \item Quyết định số 671/\gls{qd-btttt} ngày 26/04/2024 của Bộ Thông tin và Truyền thông về hướng dẫn xác định chi phí phần mềm nội bộ;
    \item Quyết định số 320/QĐ-BKHCN ngày 30/3/2025 về hướng dẫn xác định đơn giá nhân công trong quản lý chi phí đầu tư ứng dụng \gls{cntt};
    \item Nghị quyết số 36-NQ/TW ngày 01/9/2014 của Bộ Chính trị về đẩy mạnh ứng dụng và phát triển \gls{cntt} đáp ứng yêu cầu phát triển bền vững và hội nhập quốc tế;
    \item Quyết định số 749/QĐ-TTg ngày 03/6/2020 của Thủ tướng Chính phủ phê duyệt Chương trình Chuyển đổi số quốc gia đến năm 2025, định hướng đến năm 2030;
    \item Quyết định của \gls{ubnd} Thành phố Đà Nẵng về phê duyệt chủ trương đầu tư ứng dụng công nghệ \gls{gis} và \gls{ai} trong quản lý hạ tầng đô thị năm 2025.
\end{itemize}

\subsection{Tên dự án}
\textbf{\gls{geoai} - Ứng dụng \gls{gis} và Trí tuệ Nhân tạo trong Quản lý Tài sản Đô thị trên Bản đồ.}

\noindent \textbf{Phạm vi áp dụng:} Các cơ quan quản lý đô thị thuộc Thành phố Đà Nẵng, bao gồm \gls{stttt}, Sở Xây dựng, Sở Giao thông Vận tải.

\vspace{0.2cm}
\noindent \textbf{Tiến độ triển khai:} 
Đã triển khai giai đoạn 1~$\boxtimes$ \quad 
giai đoạn 2~$\square$ \quad 
giai đoạn 3~$\square$ \quad 
giai đoạn 4~$\square$ 

\subsection{Tích hợp trên hệ thống}
\textbf{Tích hợp Cổng dữ liệu đô thị TP. Đà Nẵng:} Có

\subsection{Tổng dự toán}
511,758,071 VNĐ

\subsection{Thời gian thực hiện}
1/3/2025 - 1/6/2025

\subsection{Yêu cầu, phạm vi theo chủ trương được phê duyệt}
\begin{itemize}
    \item Có kế hoạch PM, quản lý phạm vi - thay đổi.
    \item Ứng dụng \gls{webgis}/Local \gls{gis} có \gls{ai} hỗ trợ truy vấn, phân tích, dữ liệu mẫu.
\end{itemize}
\subsection{Khái quát nội dung thực hiện}
Trong phạm vi dự án năm 2025-2026, các công việc trọng tâm thực hiện bao gồm:
\begin{itemize}
    \item \textbf{Công việc 1:}  Xây dựng hệ thống GIS cốt lõi và cơ sở dữ liệu tài sản đô thị.
    \item \textbf{Công việc 2:} Phát triển module AI nhận dạng và phân tích tình trạng tài sản.
    \item \textbf{Công việc 3:} Xây dựng giao diện ứng dụng.
    \item \textbf{Công việc 4:} Kiểm thử chức năng và kiểm thử an toàn an ninh thông tin.
    \item \textbf{Công việc 5:} Nhập liệu dữ liệu tài sản ban đầu và tích hợp hệ thống.
    \item \textbf{Công việc 6:} Cài đặt, đào tạo và hướng dẫn sử dụng.
\end{itemize}
% --- BẮT ĐẦU BẢNG ---

    

% Sử dụng cột c (canh giữa) cho STT và cột p (đoạn văn có kích thước cố định) cho nội dung để chữ tự xuống dòng
\begin{longtable}{|c|p{0.85\textwidth}|}
        \caption{Danh sách yêu cầu chức năng của hệ thống GeoAI} \label{tab:yeu_cau_chuc_nang} \\ % Thêm dòng này

\hline
\rowcolor{gray!25} % Tô màu xám nhạt cho hàng tiêu đề
\textbf{STT} & \centering\arraybackslash\textbf{Yêu cầu chức năng} \\ \hline
\endfirsthead

% Định nghĩa phần đầu của các trang tiếp theo (chỉ có đường kẻ ngang, KHÔNG lặp lại tiêu đề)
\hline
\endhead

% Định nghĩa phần cuối của các trang (khi bị ngắt trang)
\hline
\endfoot

% Định nghĩa phần cuối cùng của toàn bộ bảng
\hline
\endlastfoot

% --- NỘI DUNG BẢNG ---

\textbf{I} & \textbf{Chức năng hệ thống} \\ \hline
1 & Quản lý danh mục Cấu hình hệ thống \\ \hline
2 & Quản lý danh mục \gls{api} Key \\ \hline
3 & Quản lý danh mục Nhóm quyền \\ \hline
4 & Quản lý danh mục Chức năng \\ \hline
5 & Log \gls{api} (các phần mềm khác sử dụng \gls{api} của hệ thống) \\ \hline
6 & Nhật ký sử dụng hệ thống \\ \hline
7 & Sao lưu dữ liệu \\ \hline
8 & Phân quyền cho người dùng \\ \hline

\textbf{II} & \textbf{Chức năng \gls{webgis} -- Hiển thị bản đồ} \\ \hline
9 & Hiển thị bản đồ đô thị nền (OSM / Google Maps / VN Map), zoom 1--22 \\ \hline
10 & Quản lý lớp dữ liệu: bật/tắt lớp, kéo thả thứ tự, điều chỉnh độ trong suốt \\ \hline
11 & Hiển thị tài sản trên bản đồ với icon phân loại và popup thông tin khi click \\ \hline
12 & Tìm kiếm địa chỉ, vị trí và tài sản theo mã/tên; tự động zoom đến kết quả \\ \hline
13 & Lọc hiển thị tài sản theo loại, trạng thái, quận/phường, ngày cập nhật \\ \hline
14 & Vẽ và chỉnh sửa đối tượng không gian (điểm, đường, vùng) trực tiếp trên bản đồ \\ \hline
15 & Đo khoảng cách và diện tích trực tiếp trên bản đồ \\ \hline
16 & Xuất bản đồ sang PNG/PDF kèm legend, tỷ lệ; chia sẻ bản đồ qua link URL \\ \hline

\textbf{III} & \textbf{Chức năng Quản lý Tài sản} \\ \hline
17 & Thêm, sửa, xóa tài sản với đầy đủ thông tin thuộc tính, tọa độ và hình ảnh \\ \hline
18 & Quản lý hồ sơ tài sản: lịch sử bảo trì, tài liệu kỹ thuật, trạng thái và giá trị \\ \hline
19 & Quản lý lịch bảo trì định kỳ, cảnh báo đến hạn và ghi nhận kết quả thực hiện \\ \hline
20 & Import tài sản từ Excel/CSV/Shapefile; export sang Excel/GeoJSON/Shapefile \\ \hline
21 & Dashboard tổng quan tình trạng tài sản theo loại, trạng thái, đơn vị hành chính \\ \hline

\textbf{IV} & \textbf{Chức năng \gls{ai} -- Truy vấn thông minh} \\ \hline
22 & Truy vấn tài sản bằng ngôn ngữ tự nhiên tiếng Việt, trả kết quả trên bản đồ và bảng dữ liệu \\ \hline
23 & Hiển thị và cho phép chỉnh sửa câu SQL được sinh ra từ câu hỏi tự nhiên \\ \hline
24 & \gls{ai} đề xuất kế hoạch bảo trì dự đoán dựa trên lịch sử và tuổi thọ tài sản \\ \hline
25 & Nhận dạng tình trạng tài sản từ ảnh chụp hiện trường (hỏng hóc, vết nứt, gỉ sét) \\ \hline
26 & Tự động sinh báo cáo tình trạng tài sản bằng ngôn ngữ tự nhiên tiếng Việt, xuất PDF \\ \hline

\textbf{V} & \textbf{Chức năng Phân tích Không gian} \\ \hline
27 & Phân tích mật độ và phân bố tài sản theo không gian (heatmap) \\ \hline
28 & Phân tích vùng đệm (buffer) xung quanh tài sản, thống kê tài sản trong vùng \\ \hline
29 & Tối ưu tuyến đường bảo trì, tính lộ trình ngắn nhất cho danh sách tài sản cần xử lý \\ \hline
30 & Thống kê tài sản theo đơn vị hành chính: phường, quận, thành phố; bản đồ choropleth \\ 
\end{longtable}

\subsection{Tuân thủ khung kiến trúc}
Dự án \gls{geoai} được xây dựng tuân thủ Kiến trúc Chính phủ điện tử Việt Nam phiên bản 3.0, Khung kiến trúc đô thị thông minh của Bộ Thông tin và Truyền thông, và chuẩn OGC (Open Geospatial Consortium) trong lĩnh vực \gls{gis}.

\subsection{Đề xuất cấp độ an toàn thông tin}
Hệ thống \gls{geoai} đề xuất cấp độ an toàn hệ thống thông tin Cấp 3 theo Nghị định số 85/2016/\gls{nd-cp}, do hệ thống quản lý dữ liệu quan trọng về hạ tầng đô thị và có kết nối với các hệ thống cốt lõi của thành phố.

\cleardoublepage